\scanpage{39}
\mspage{- 29 -}
% archive sheet number: 38

\[
\begin{aligned}
\frac{\handwritten{\partial}}{\handwritten{\partial}x_k}\left(\frac{\handwritten{\partial u}_j}{\handwritten{\partial}x_i}\right)&=\frac{\handwritten{\partial}U_{ij}}{\handwritten{\partial}x_k}(\handwritten{x};\handwritten{u}(\handwritten{x}))+\handwritten{\sum_{\unsure{k|h}=1}^q}\frac{\handwritten{\partial}U_{ij}}{\handwritten{\partial}y_h}(\handwritten{x};\handwritten{u}(\handwritten{x}))\mathbin{.}\frac{\handwritten{\partial u}_h}{\handwritten{\partial}x_k}
&& (1\handwritten{\leq}i\handwritten{\leq}p,\\
&=\frac{\handwritten{\partial}U_{ij}}{\handwritten{\partial}x_k}(\handwritten{x};\handwritten{u}(\handwritten{x}))+\handwritten{\sum_{\unsure{k|h}=1}^q}\frac{\handwritten{\partial}U_{ij}}{\handwritten{\partial}y_h}(\handwritten{x};\handwritten{u}(\handwritten{x}))U_{kh}(\handwritten{x};\handwritten{u}(\handwritten{x})).
&&1\handwritten{\leq}k\handwritten{\leq}p,\\
&&&1\handwritten{\leq}j\handwritten{\leq}q
\end{aligned}
\]

Ecriva\handwritten{-}nt que les expressions ainsi obtenues pour $\frac{\handwritten{\partial}}{\handwritten{\partial}x_k}(\frac{\handwritten{\partial u}_j}{\handwritten{\partial}x_i})$ et $\frac{\handwritten{\partial}}{\handwritten{\partial}x_i}(\frac{\handwritten{\partial u}_j}{\handwritten{\partial}x_k})$ sont égales au point $\handwritten{x}$, puis faisant $\handwritten{x}=\handwritten{x}_0$ et tenant compte de (23) il vient
\[
\begin{aligned}
&\frac{\handwritten{\partial}U_{ij}}{\handwritten{\partial}x_k}(\handwritten{x}_0;\handwritten{y}_0)+\handwritten{\sum_{h=1}^q}\frac{\handwritten{\partial}U_{ij}}{\handwritten{\partial}y_h}(\handwritten{x}_0;\handwritten{y}_0)U_{kh}(\handwritten{x}_0;\handwritten{y}_0)=\\
&=\frac{\handwritten{\partial}U_{kj}}{\handwritten{\partial}x_i}(\handwritten{x}_0;\handwritten{y}_0)+\handwritten{\sum_{\unsure{k|h}=1}^q}\frac{\handwritten{\partial}U_{kj}}{\handwritten{\partial}y_h}(\handwritten{x}_0;\handwritten{y}_0)U_{jh}(\handwritten{x}_0;\handwritten{y}_0)
\end{aligned}
\]
Mais comme par hypothèse, $(\handwritten{x}_0,\handwritten{y}_0)$ est un point arbitraire de A, on obtient bien (24).

$2^{\handwritten{\circ}}$ Les conditions sont \uline{suffisantes}. Considérons l'équation différentielle
\[
\handwritten{y}'=\underline{U}(\handwritten{x}_0+t(\handwritten{x}-\handwritten{x}_0);\handwritten{y})\mathbin{.}(\handwritten{x}-\handwritten{x}_0)\tag*{(25)}
\]

où $\underline{U}$ est la matrice $(U_{ij})$, $\handwritten{x}_0$ et $\handwritten{x}$ sont fixés dans B, et $\handwritten{y}'=d\handwritten{y}/dt$, t variant dans un voisina ge de 0 dans K. Il résulte de la formation de la série majorante dans la démonstration du th.2 que, dès que $\handwritten{x}-\handwritten{x}_0$ est assez petit, l'équation (25) admet une solution t une seule $\handwritten{v}(t;\handwritten{x}_0,\handwritten{x},\handwritten{y}_0)$ holomorphe pour $\handwritten{|}t\handwritten{|<}2$ et pour $\handwritten{x}$ et $\handwritten{x}_0$ dans un voisinage assez petit U de $\handwritten{a}$, $\handwritten{y}_0$ dans un voisinage assez petit V de $\handwritten{b}$ et telle que $\handwritten{v}(0;\handwritten{x}_0,\handwritten{x},\handwritten{y}_0)=\handwritten{y}_0$. Si $\handwritten{v}=(v_j)$, on a donc
\[
\frac{dv_j}{dt}=\handwritten{\sum_{i=1}^p}U_{ij}(\handwritten{x}_0+t(\handwritten{x}-\handwritten{x}_0);\handwritten{v}(t))(x_i-x_{i0})\qquad(1\handwritten{\leq}j\handwritten{\leq}q)\tag*{(26)}
\]
On déduit de là que
\[
\begin{aligned}
\frac{d}{dt}\left(\frac{\handwritten{\partial}v_j}{\handwritten{\partial}x_k}\right)&=\frac{\handwritten{\partial}}{\handwritten{\partial}x_k}\left(\frac{dv_j}{dt}\right)=U_{kj}(\handwritten{x}_0+t(\handwritten{x}-\handwritten{x}_0);\handwritten{v}(t))+\\
&\quad+t\handwritten{\sum_i}(x_i-x_{i0})\frac{\handwritten{\partial}U_{ij}}{\handwritten{\partial}x_k}(\handwritten{x}_0+t(\handwritten{x}-\handwritten{x}_0);\handwritten{v}(t))\\
&\quad+\handwritten{\sum_{i,h}}(x_i-x_{i0})\frac{\handwritten{\partial}U_{ij}}{\handwritten{\partial}y_h}(\handwritten{x}_0+t(\handwritten{x}-\handwritten{x}_0);\handwritten{v}(t))\frac{\handwritten{\partial}v_{\unsure{n|h}}}{\handwritten{\partial}x_k}
\qquad(1\handwritten{\leq}j\handwritten{\leq}q,\ 1\handwritten{\leq}k\handwritten{\leq}p)
\end{aligned}
\]
